\documentclass[11pt,a4paper]{article}

% ============================================================
% PACOTES
% ============================================================

\usepackage[T1]{fontenc}
\usepackage[utf8]{inputenc}
\usepackage[brazil]{babel}

\usepackage{lmodern}
\usepackage{microtype}

\usepackage{amsmath,amssymb,amsthm,mathtools}
\usepackage{enumitem}

\usepackage{geometry}
\usepackage{setspace}

\usepackage{xcolor}
\usepackage{hyperref}
\usepackage{cleveref}

% ============================================================
% FORMATAÇÃO
% ============================================================

\geometry{
    top=2.7cm,
    bottom=2.7cm,
    left=2.7cm,
    right=2.7cm
}

\onehalfspacing

\setlength{\parindent}{1.25cm}
\setlength{\parskip}{0pt}

\hypersetup{
    colorlinks=true,
    linkcolor=blue!50!black,
    citecolor=blue!50!black,
    urlcolor=blue!50!black
}

% ============================================================
% AMBIENTES
% ============================================================

\newtheorem{theorem}{Teorema}[section]
\newtheorem{proposition}[theorem]{Proposição}
\newtheorem{lemma}[theorem]{Lema}
\newtheorem{corollary}[theorem]{Corolário}
\newtheorem{conjecture}[theorem]{Conjectura}

\theoremstyle{definition}
\newtheorem{definition}[theorem]{Definição}
\newtheorem{question}[theorem]{Questão}
\newtheorem{problem}[theorem]{Problema}

\theoremstyle{remark}
\newtheorem{remark}[theorem]{Observação}

% ============================================================
% COMANDOS
% ============================================================

\newcommand{\N}{\mathbb{N}}
\newcommand{\Z}{\mathbb{Z}}
\newcommand{\R}{\mathbb{R}}
\newcommand{\cS}{\mathcal{S}}
\newcommand{\cZ}{\mathcal{Z}}
\newcommand{\abs}[1]{\left|#1\right|}

% ============================================================
% DOCUMENTO
% ============================================================

\begin{document}

% ============================================================
% CABEÇALHO
% ============================================================

\begin{center}

  {\LARGE\bfseries
    Pré-projeto
  }

\end{center}

\vspace{0.8cm}

% ============================================================
% RESUMO
% ============================================================

\noindent\textbf{Resumo.}

Este projeto investiga problemas estruturais, extremais e enumerativos
para famílias de conjuntos de inteiros nas quais se exige que
representações aditivas distintas permaneçam \emph{quantitativamente
  separadas}. O protótipo é a classe dos conjuntos Sidon \(k\)-fortes,
para a qual já obtivemos, durante o doutorado, um teorema estrutural, a
determinação assintótica do problema extremal e limites superior e
inferior para o problema de contagem com a mesma dependência em \(n\) e
\(k\).

O resultado estrutural mostra que a condição \(k\)-forte, aparentemente
uma restrição apenas sobre certas configurações ordenadas de quatro
elementos, força a separação dos próprios elementos: em um conjunto
Sidon \(k\)-forte há no máximo um par a distância menor que \(k\). Isso
estabelece uma ponte com os \(B_{2,\Delta}\)-sets recentemente
introduzidos por Nathanson, e sugere um programa mais amplo envolvendo
\(B_h\)-sets, \(B_{h,\Delta}\)-sets e o método de contêineres.

Os objetivos centrais são determinar as constantes corretas nos
problemas de contagem --- inclusive no caso clássico \(k=1\), onde o
problema de Cameron--Erdős permanece em aberto ---, entender o alcance
do fenômeno estrutural descoberto, e desenvolver a teoria correspondente
para \(h>2\).

\vspace{0.4cm}
\tableofcontents
\newpage

% ============================================================
\section{Introdução}
% ============================================================

\subsection{Conjuntos de Sidon}

Um conjunto \(A\) de inteiros é um \emph{conjunto de Sidon} se todas as
somas \(a+b\) com \(a,b\in A\) e \(a\leq b\) são distintas;
equivalentemente, se a equação \(a+b=c+d\) admite apenas as soluções
triviais, aquelas com \(\{a,b\}=\{c,d\}\).

Escrevendo \([n]:=\{1,\dots,n\}\) e
\[
  F(n):=\max\{|A|:A\subseteq[n]\text{ é Sidon}\},
\]
resultados clássicos de Erdős, Turán, Chowla e Singer, dos anos 1940,
estabelecem \(F(n)=(1+o(1))\sqrt n\). O problema extremal está, portanto,
resolvido em ordem principal.

Uma generalização natural são os \(B_h\)-sets: \(A\) é \(B_h\) quando
cada elemento do conjunto-soma \(hA=A+\dots+A\) admite uma única
representação como soma de \(h\) elementos de \(A\) tomados em ordem não
decrescente. Para \(h=2\) recuperamos os conjuntos de Sidon.

\subsection{O problema enumerativo}

Além de perguntar qual o maior conjunto de Sidon em \([n]\), é natural
perguntar \emph{quantos} existem. Escrevendo
\[
  A(n):=\bigl|\{A\subseteq[n]:A\text{ é Sidon}\}\bigr|,
\]
uma cota inferior trivial é \(A(n)\geq2^{F(n)}=2^{(1+o(1))\sqrt n}\),
pois todo subconjunto de um conjunto de Sidon é de Sidon.

Para muitas famílias de estruturas livres de subestruturas proibidas
vale a heurística
\begin{equation}
  \label{eq:heuristica}
  |\{\text{estruturas}\}|
  =
  2^{(1+o(1))\cdot\text{número extremal}},
\end{equation}
isto é, a cota trivial é essencialmente correta. Cameron e Erdős
conjecturaram, nos anos 1980, que~\eqref{eq:heuristica} valeria para
conjuntos de Sidon. A conjectura é falsa: Saxton e Thomason exibiram uma
construção fornecendo
\[
  A(n)\geq2^{(\frac12\log_25+o(1))\sqrt n}
  =
  2^{(1{,}1609\ldots+o(1))\sqrt n},
\]
e Kohayakawa, Lee, Rödl e Samotij provaram
\(A(n)\leq2^{(6{,}442+o(1))\sqrt n}\). A determinação da constante
correta permanece um problema em aberto, e o intervalo entre
\(1{,}16\) e \(6{,}44\) é um dos poucos casos em que se sabe
que~\eqref{eq:heuristica} falha sem que se saiba por quanto.

\subsection{Conjuntos Sidon \texorpdfstring{\(k\)}{k}-fortes}

Uma reformulação útil da propriedade de Sidon é a seguinte: \(A\) é
Sidon se e somente se
\[
  |(x+w)-(y+z)|\geq1
  \qquad\text{sempre que}\qquad
  x<y\leq z<w,\quad x,y,z,w\in A.
\]
Isso sugere impor uma separação maior.

\begin{definition}
  \label{def:k-forte}
  Um conjunto \(A\subseteq\mathbb Z\) é \emph{Sidon \(k\)-forte} se
  \[
    |(x+w)-(y+z)|\geq k
    \qquad\text{para todos}\qquad
    x<y\leq z<w\ \text{em}\ A .
  \]
\end{definition}

Para \(k=1\) recuperamos a propriedade de Sidon usual. Enquanto a
condição clássica proíbe colisões exatas entre certas somas, a
propriedade \(k\)-forte exige que essas somas permaneçam a uma distância
mínima prescrita. Essa noção foi introduzida por Kohayakawa, Lee,
Moreira e Rödl no contexto de conjuntos de Sidon infinitos contidos em
conjuntos aleatórios de inteiros, onde consideram o caso \(k=n^\alpha\).

Uma noção aparentada, mais forte, foi introduzida recentemente por
Nathanson.

\begin{definition}
  \label{def:bhdelta}
  Para \(A\subseteq\mathbb Z\), denotamos por \(A\oplus A\) o
  \emph{multiconjunto} das somas \(a+b\) com \(a\leq b\) em \(A\), no
  qual representações distintas de um mesmo inteiro aparecem como
  ocorrências distintas. Dizemos que um multiconjunto de inteiros é
  \emph{\(\Delta\)-separado} se quaisquer duas ocorrências distintas
  diferem em pelo menos \(\Delta\). Um \emph{\(B_{2,\Delta}\)-set} é um
  conjunto \(A\) tal que \(A\oplus A\) é \(\Delta\)-separado.
\end{definition}

À primeira vista as duas definições são bastante distintas: a
Definição~\ref{def:k-forte} restringe apenas configurações ordenadas
específicas, enquanto a Definição~\ref{def:bhdelta} impõe separação a
\emph{todas} as somas. Um dos resultados já obtidos mostra que elas
estão, de fato, a um único elemento de distância.

% ============================================================
\section{Resultados obtidos}
% ============================================================

Os resultados abaixo, desenvolvidos durante o doutorado, constituem o
ponto de partida do projeto. Escrevemos
\[
  F_k(n):=\max\{|A|:A\subseteq[n]\text{ é Sidon \(k\)-forte}\},
  \qquad
  N_k(n):=\bigl|\{A\subseteq[n]:A\text{ é Sidon \(k\)-forte}\}\bigr| .
\]

\begin{center}
  \renewcommand{\arraystretch}{1.4}
  \begin{tabular}{llc}
    \hline
    \textbf{Resultado}              & \textbf{Enunciado}                                & \textbf{Seção}     \\
    \hline
    Estrutural                      & núcleo \(k\)-separado após remover \(1\) elemento
                                    & \ref{sec:estrutural}                                                   \\
    Caracterização                  & \(k\)-forte \(+\) \(k\)-separado
    \(\Leftrightarrow\) \(B_{2,k}\) & \ref{sec:b2k}                                                          \\
    Extremal                        & \(F_k(n)=(1+o(1))\sqrt{n/k}\)                     & \ref{sec:extremal} \\
    Contagem (inferior)             & \(N_k(n)\geq2^{c\log_2(k+1)\sqrt{n/k}}\)
                                    & \ref{sec:lower}                                                        \\
    Contagem (superior)             & \(N_k(n)\leq2^{C\log_2(k+1)\sqrt{n/k}}\)
                                    & \ref{sec:upper}                                                        \\
    \hline
  \end{tabular}
\end{center}

\subsection{Um teorema estrutural de separação}
\label{sec:estrutural}

O ponto de partida é a seguinte observação, que chamamos de
\emph{propriedade de quatro pontos}: se \(x_1<x_2<x_3<x_4\) pertencem a
um conjunto Sidon \(k\)-forte, então \emph{qualquer} partição desses
quatro elementos em dois pares produz somas cuja diferença tem módulo ao
menos \(k\) --- e não apenas a partição \(\{x_1,x_4\}\mid\{x_2,x_3\}\),
que é a única diretamente contemplada pela Definição~\ref{def:k-forte}.

A razão é simples: escrevendo \(p=x_2-x_1\), \(q=x_3-x_2\) e
\(r=x_4-x_3\), as três diferenças de somas de pares são
\[
  |r-p|,\qquad p+r,\qquad p+2q+r,
\]
e a primeira, que é a coberta pela definição, é a menor das três.

Essa observação tem uma consequência estrutural forte.

\begin{theorem}[Núcleo separado]
  \label{thm:separated-core}
  Seja \(A\subseteq\mathbb Z\) um conjunto Sidon \(k\)-forte. Então
  existe no máximo um par de elementos distintos \(a,b\in A\) com
  \(|a-b|<k\). Consequentemente, existe \(A'\subseteq A\) com
  \(|A'|\geq|A|-1\) tal que \(|a-b|\geq k\) para todos os elementos
  distintos de \(A'\).
\end{theorem}

A prova considera o grafo cujas arestas registram os pares a distância
menor que \(k\). Dois pares curtos disjuntos produziriam quatro
elementos violando a propriedade de quatro pontos; dois pares curtos
compartilhando um elemento produziriam \(x<y<z\) com \(|x+z-2y|<k\),
contradizendo diretamente a Definição~\ref{def:k-forte} aplicada a
\(x<y\leq y<z\). Segue que o grafo é uma estrela, e basta remover seu
centro.

O enunciado é mais forte do que a mera existência de um subconjunto
separado grande: toda a falha de separação está concentrada em um único
par.

\subsection{Separação de somas e \texorpdfstring{\(B_{2,k}\)}{B2k}-sets}
\label{sec:b2k}

\begin{proposition}
  \label{prop:b2k}
  Para \(A\subseteq\mathbb Z\), o multiconjunto \(A\oplus A\) é
  \(k\)-separado se, e somente se, \(A\) é Sidon \(k\)-forte e
  \(k\)-separado.
\end{proposition}

A implicação direta é imediata: a separação de \(A\oplus A\) fornece a
condição \(k\)-forte comparando \(x+w\) com \(y+z\), e a separação de
\(A\) comparando \(a+a\) com \(a+b\). Na recíproca, a propriedade de
quatro pontos controla os pares com quatro elementos distintos,
enquanto a \(k\)-separação de \(A\) controla os casos em que as duas
somas compartilham algum elemento.

Combinando com o Teorema~\ref{thm:separated-core}, obtemos:

\begin{corollary}
  Todo conjunto Sidon \(k\)-forte contém, após a remoção de no máximo um
  elemento, um \(B_{2,k}\)-set.
\end{corollary}

Essa é a ponte concreta entre a classe estudada no doutorado e os
objetos introduzidos por Nathanson, e é ela que motiva boa parte dos
objetivos deste projeto.

\subsection{O problema extremal}
\label{sec:extremal}

A cota inferior vem de uma dilatação: se \(A\) é Sidon, então
\(kA:=\{ka:a\in A\}\) é Sidon \(k\)-forte, pois todas as diferenças de
somas ficam multiplicadas por \(k\). Aplicando isso a um conjunto de
Sidon máximo em \([\lfloor n/k\rfloor]\), obtemos
\(F_k(n)\geq(1-o(1))\sqrt{n/k}\).

Na direção oposta, o Teorema~\ref{thm:separated-core} permite reduzir o
problema a um conjunto \(k\)-separado, ao custo de um elemento. Um
argumento de dupla contagem com intervalos deslizantes, no espírito de
Erdős--Turán, fornece então
\[
  F_k(n)
  \leq
  \sqrt{\frac nk}
  +
  O\!\left(\left(\frac nk\right)^{1/3}\right),
\]
e portanto \(F_k(n)=(1+o(1))\sqrt{n/k}\) sempre que \(k=o(n)\).

Além de resolver o problema extremal em ordem principal, essa estimativa
é ingrediente essencial do problema enumerativo: ela restringe a faixa
de cardinalidades sobre a qual é preciso somar na contagem.

\subsection{Contagem: o limite inferior}
\label{sec:lower}

A construção generaliza a de Saxton e Thomason em duas camadas.
Partimos de um conjunto de Sidon \(A_0\) num grupo cíclico
\(\mathbb Z_m\) e, para cada elemento, escolhemos independentemente uma
entre \(1+\ell(\beta+1)\) alternativas: omiti-lo, ou colocá-lo em um
dentre \(\ell\) blocos consecutivos e atribuir-lhe uma dentre
\(\beta+1\) perturbações. Em seguida dilatamos por um fator \(\alpha\).

A primeira camada preserva a propriedade de Sidon; a segunda a converte
na propriedade \(k\)-forte, desde que \(\alpha-2\beta\geq k\) --- a
dilatação separa as somas por \(\alpha\), e as perturbações reduzem essa
separação em no máximo \(2\beta\). Como a construção é injetiva, cada
escolha produz um conjunto distinto, e o custo espacial é governado por
\(\alpha\ell\). O resultado é
\[
  \log_2 N_k(n)
  \geq
  (1-o(1))\,
  \frac{\log_2\bigl(1+\ell(\beta+1)\bigr)}{\sqrt{\alpha\ell}}\,\sqrt n .
\]

Escolhas simples já fornecem uma constante absoluta. Tomando
\((\alpha,\beta,\ell)=(3k,k,1)\), obtemos:

\begin{theorem}
  \label{thm:lower-count}
  Para cada inteiro \(k\geq1\) fixo e todo \(n\) suficientemente grande
  (com \(n_0\) dependendo de \(k\)),
  \[
    N_k(n)\geq2^{\,\frac1{\sqrt3}\log_2(k+1)\sqrt{n/k}} .
  \]
\end{theorem}

Essa escolha não é ótima. Para \(k=1\), a escolha
\((\alpha,\beta,\ell)=(1,0,4)\) recupera exatamente a construção de
Saxton--Thomason, com constante \(\tfrac12\log_25=1{,}1609\ldots\); para
\(k=2\), a escolha \((4,1,2)\) mantém as mesmas cinco alternativas por
elemento --- dois blocos, duas perturbações em cada, mais a omissão ---
fornecendo constante \(\log_25/\sqrt8=0{,}8209\ldots\). Otimizando os
três parâmetros quando \(k\to\infty\), a escala natural é
\(\beta\asymp k/\log k\), e a constante normalizada tende a \(1\).

\subsection{Contagem: o limite superior}
\label{sec:upper}

O limite superior usa o método de \emph{contêineres de grafos}. A ideia
geral do método é a seguinte: em um grafo suficientemente denso, todo
conjunto independente fica contido na união de um pequeno
\emph{certificado} com um \emph{contêiner} pequeno; como há poucos
certificados e poucos contêineres, e cada par determina o conjunto, o
número de independentes fica controlado.

Para aplicá-lo aqui, fixa-se um conjunto parcial \(U\) e constrói-se um
multigrafo auxiliar cujos vértices são os candidatos a novos elementos e
cujas arestas registram quase-colisões com elementos de \(U\). O
Teorema~\ref{thm:separated-core} permite substituir \(U\), com perda de
um elemento, por um subconjunto \(k\)-separado --- e essa separação tem
uma consequência decisiva: a multiplicidade das arestas do multigrafo
passa a ser no máximo \(2\), o que permite passar para o grafo simples
subjacente perdendo apenas um fator constante.

Um resultado de \emph{supersaturação} --- uma cota inferior para o
número de arestas induzidas por qualquer subconjunto grande de vértices
--- garante então a densidade necessária, e a iteração do lema de
contêineres fornece uma estimativa para conjuntos de cardinalidade
prescrita. Somando sobre as cardinalidades possíveis, e usando a
estimativa extremal da Seção~\ref{sec:extremal} para limitar essa soma,
obtemos:

\begin{theorem}
  \label{thm:upper-count}
  Existe uma constante absoluta \(C>0\) tal que, para \(n\)
  suficientemente grande e \emph{uniformemente} em
  \(1\leq k\leq n/(\log n)^2\),
  \[
    N_k(n)\leq2^{\,C\log_2(k+1)\sqrt{n/k}} .
  \]
\end{theorem}

Combinando os Teoremas~\ref{thm:lower-count} e~\ref{thm:upper-count},
para cada \(k\) fixo,
\[
  \log_2 N_k(n)\asymp\log(k+1)\sqrt{\frac nk} .
\]
Convém registrar a diferença de uniformidade entre os dois enunciados: o
limite superior é uniforme na faixa indicada, enquanto no limite
inferior o valor de \(n_0\) depende de \(k\). Tornar o limite inferior
uniforme, e determinar a constante correta, são objetivos do projeto.

Observamos ainda que, em vista do resultado extremal
\(F_k(n)=(1+o(1))\sqrt{n/k}\), a presença do fator \(\log(k+1)\) mostra
que a heurística~\eqref{eq:heuristica} falha para \emph{toda} a família,
e não apenas no caso clássico \(k=1\).

% ============================================================
\section{Objetivos e questões de pesquisa}
% ============================================================

O objetivo geral é desenvolver uma teoria estrutural e enumerativa para
famílias aditivas nas quais representações distintas devem permanecer
quantitativamente separadas. Os objetivos abaixo estão ordenados por
grau de maturidade: os dois primeiros consolidam e refinam resultados já
obtidos; os demais são progressivamente mais exploratórios.

\subsection{Objetivo I --- Constantes na contagem de conjuntos de Sidon}

Como discutido na Introdução, a constante correta em
\(A(n)=2^{\Theta(\sqrt n)}\) é desconhecida, com
\(1{,}16\leq c\leq6{,}44\). A pergunta é se os argumentos de
supersaturação e contêineres desenvolvidos para o caso \(k\)-forte
podem, especializados a \(k=1\), reduzir esse intervalo.

Duas observações motivam a tentativa. Primeiro, o argumento
multifásico usual perde constantes em vários pontos independentes
(escolha dos conjuntos-teste, tamanho do contêiner final, estimativas
binomiais); um objetivo concreto é \emph{quantificar} onde está a perda
dominante. Segundo, no caso \(k\)-forte com \(k\to\infty\) as constantes
das estimativas intermediárias entram aditivamente dentro de um
logaritmo e são absorvidas pelo fator \(\log k\) --- entender por que
isso não ocorre em \(k=1\) pode indicar onde está o obstáculo real.

\subsection{Objetivo II --- A constante no problema \texorpdfstring{\(k\)}{k}-forte}

Os resultados obtidos determinam a escala \(\log(k+1)\sqrt{n/k}\).
Pretende-se determinar a constante, isto é, estudar
\[
  c_k
  :=
  \lim_{n\to\infty}
  \frac{\log_2 N_k(n)}{\log_2(k+1)\sqrt{n/k}}
\]
--- começando por estabelecer que esse limite existe --- e o seu
comportamento quando \(k\to\infty\).

Há indícios de que a resposta seja \(c_k\to1\). Do lado inferior, a
otimização dos parâmetros fornece constante \((1+o(1))\log_2k/\sqrt k\).
Do lado superior, o termo dominante da estimativa é da forma
\((Ck)^{F_k(n)}\) para uma constante absoluta \(C\), e
\(\log_2(Ck)/\log_2k\to1\). Se ambas as pontas puderem ser tornadas
precisas, obteríamos
\[
  N_k(n)=2^{(1+o(1))\log_2k\sqrt{n/k}},
\]
um resultado sem análogo conhecido no caso \(k=1\).

Questões associadas:
\begin{itemize}
  \item tornar o limite inferior uniforme em \(k\);
  \item ampliar a faixa \(k\leq n/(\log n)^2\) do limite superior;
  \item obter estimativas para conjuntos de cardinalidade prescrita.
\end{itemize}

\subsection{Objetivo III --- Estrutura de conjuntos \texorpdfstring{\(k\)}{k}-fortes}

O Teorema~\ref{thm:separated-core} sugere uma rigidez estrutural ainda
não explorada. Que outras propriedades de um \(B_{2,k}\)-set
sobrevivem, à custa de poucos elementos, em conjuntos Sidon
\(k\)-fortes? E o fenômeno ``remover poucos elementos produz separação
completa'' tem análogos para outras equações aditivas, como as que
definem conjuntos livres de somas ou livres de progressões?

\subsection{Objetivo IV --- \texorpdfstring{\(B_h\)}{Bh}-sets e contêineres de hipergrafos}

Para a contagem de \(B_h\)-sets de cardinalidade prescrita, Dellamonica,
Kohayakawa, Lee, Rödl e Samotij obtiveram estimativas precisas por
métodos ad hoc, combinando supersaturação, grafos de colisão e
conjuntos-semente.

Pretende-se revisitar o problema com a linguagem de contêineres de
hipergrafos. O hipergrafo natural tem vértices \([n]\) e hiperarestas
associadas às soluções não triviais de
\(a_1+\dots+a_h=b_1+\dots+b_h\); os \(B_h\)-sets são seus conjuntos
independentes. Uma aplicação direta dos teoremas de contêineres fornece
informação enumerativa, mas não recupera as melhores estimativas
conhecidas. A questão metodológica é identificar exatamente qual
informação adicional --- sobre supersaturação, graus ou codegrees --- é
necessária para fechar essa diferença.

\subsection{Objetivo V --- \texorpdfstring{\(B_{h,\Delta}\)}{Bh,Delta}-sets}

Nathanson introduziu os \(B_{h,\Delta}\)-sets exigindo que somas
distintas de \(h\) elementos permaneçam a distância pelo menos
\(\Delta\), e propôs explicitamente o problema de estender aos
\(B_{h,\Delta}\) os resultados conhecidos para \(B_h\). Para \(h=2\), o
Corolário da Seção~\ref{sec:b2k} fornece uma ponte imediata com os
resultados já obtidos.

No problema enumerativo, a abordagem natural é substituir as colisões
exatas pelas quase-colisões
\[
  \bigl|(a_1+\dots+a_h)-(b_1+\dots+b_h)\bigr|<\Delta .
\]
Para \(\Delta\) fixo isso corresponde a uma família finita de equações
aditivas próximas, o que torna plausível a aplicação das mesmas
ferramentas. O primeiro passo é entender como a substituição afeta a
supersaturação e as estimativas de graus.

\subsection{Objetivo VI --- Uma noção \texorpdfstring{\(k\)}{k}-forte para \texorpdfstring{\(h>2\)}{h>2}}

Falta uma definição adequada de propriedade \(k\)-forte para
\(B_h\)-sets com \(h>2\). A definição desejada deve ser estritamente
mais fraca que a separação completa de \(hA\) --- para admitir uma
classe rica --- mas forçar essa separação após a remoção de poucos
elementos, como ocorre em \(h=2\).

\begin{question}
  Para cada \(h\geq2\), existe uma constante \(r(h)\) tal que todo
  conjunto satisfazendo uma noção adequada de \(k\)-forte para
  \(B_h\) contém \(A'\subseteq A\) com \(|A'|\geq|A|-r(h)\) e \(hA'\)
  \(k\)-separado?
\end{question}

O Teorema~\ref{thm:separated-core} mostra que \(r(2)=1\).

% ============================================================
\section{Metodologia}
% ============================================================

O projeto é de natureza teórica e utilizará ferramentas de Combinatória
Extremal, Aditiva e Probabilística: contêineres de grafos e hipergrafos;
resultados de supersaturação; análise de graus e codegrees; grafos e
multigrafos de colisão; funções de representação de somas e diferenças;
argumentos estruturais de separação; construções em grupos cíclicos;
dilatações, blocos e perturbações; métodos probabilísticos e argumentos
de deleção.

O mecanismo já desenvolvido para o caso \(k\)-forte --- extração de um
núcleo separado, controle de multiplicidades, supersaturação e
contêineres iterados --- fornece um roteiro. A questão metodológica
central do projeto é determinar quais de suas partes se transportam para
\(B_h\)-sets e \(B_{h,\Delta}\)-sets.

Para o Objetivo IV, o primeiro passo é quantificar a diferença entre a
aplicação direta de contêineres ao hipergrafo natural e as estimativas
finas conhecidas. Para o Objetivo V, é entender como a passagem de
colisões exatas a quase-colisões afeta supersaturação e codegrees.

% ============================================================
\section{Cronograma e organização}
% ============================================================

O projeto combina problemas em estágios distintos de maturidade, o que
permite que seu desenvolvimento não dependa da resolução de uma única
questão em aberto.

\begin{center}
  \renewcommand{\arraystretch}{1.4}
  \begin{tabular}{clc}
    \hline
    \textbf{Período}    & \textbf{Atividade}                             & \textbf{Objetivos} \\
    \hline
    Meses 1--6          & Consolidação dos resultados \(k\)-fortes;
    redação e submissão & I, II                                                               \\
    Meses 7--12         & Uniformidade e constantes na contagem          & II                 \\
    Meses 13--18        & Estrutura; transferência para \(B_{2,\Delta}\)
                        & III, V                                                              \\
    Meses 19--24        & \(B_h\)-sets via contêineres de hipergrafos    & IV                 \\
    Meses 25--30        & \(B_{h,\Delta}\)-sets e noção \(k\)-forte para
    \(h>2\)             & V, VI                                                               \\
    Meses 31--36        & Redação final e integração dos resultados      & ---                \\
    \hline
  \end{tabular}
\end{center}

% ============================================================
\section{Perspectiva geral}
% ============================================================

As diferentes linhas do projeto organizam-se em torno de uma única
questão:

\begin{quote}
  \emph{Como condições quantitativas de separação entre representações
    aditivas alteram a estrutura, o tamanho máximo e a quantidade de
    conjuntos de Sidon e \(B_h\)-sets?}
\end{quote}

No problema clássico proíbem-se colisões exatas entre somas. Nos
conjuntos Sidon \(k\)-fortes exige-se separação quantitativa apenas para
certas configurações ordenadas. Nos \(B_{h,\Delta}\)-sets a separação é
imposta a todas as representações.

O Teorema~\ref{thm:separated-core} mostra que, para \(h=2\), as duas
últimas perspectivas estão a um único elemento de distância. Do ponto de
vista extremal isso leva à escala \(\sqrt{n/k}\); do ponto de vista
enumerativo, à escala maior \(\log(k+1)\sqrt{n/k}\).

Esse fator adicional é o ponto conceitual do trabalho: a separação não
altera apenas o tamanho máximo dos conjuntos, mas a entropia da família.
Em termos informais, ao escolher um conjunto Sidon \(k\)-forte típico,
cada elemento contribui com \(\Theta(k/\log k)\) alternativas em vez de
\(\Theta(1)\) --- e é isso que produz o \(\log(k+1)\) no expoente.


\bibliographystyle{plain}
\bibliography{references}

\end{document}
